\documentclass[11pt]{article}
\usepackage[T1]{fontenc}
\usepackage{textcomp}
\usepackage[margin=0.75in]{geometry}
\usepackage{amsmath,amssymb,amsthm}
\usepackage{array,booktabs}
\usepackage{hyperref}
\setlength{\parindent}{0pt}
\newtheorem{theorem}{Theorem}
\theoremstyle{definition}
\newtheorem{definition}{Definition}
\title{Flow Proof Artifact: Order-join-trans-meet}
\author{Generated by \texttt{flow doc proof}}
\date{Flow Proof Book}
\begin{document}
\maketitle
Join is above meet via transitivity through the right argument.\\[0.5em]
\textbf{Source.} davey-priestley — *Introduction to Lattices and Order*\\[0.5em]
\bigskip
\noindent{\large \textbf{Derived fact 1.} \textit{meet(a, b) <= join(a, b)} \hfill Order \cdot lattice \cdot join is above meet via transitivity}\par
\smallskip\noindent\textit{Coordinate: Order \cdot lattice \cdot join is above meet via transitivity}\par
\smallskip\noindent\textit{Source: davey-priestley}\par
\smallskip\noindent\textit{Built on: meet is below the right argument, for meet on Order, join is above the right argument, for join on Order, less-or-equal is transitive, for order on the natural numbers}\par
\medskip
\noindent\textbf{Goal.} meet(a, b) <= join(a, b)\par
\noindent$\displaystyle \forall a \in \mathbb{N} \forall b \in \mathbb{N}\quad meet(a, b) \le join(a, b)$\par
\medskip
\noindent
\begin{tabular}{@{} >{\raggedright\arraybackslash}p{0.06\textwidth} >{\raggedright\arraybackslash}p{0.47\textwidth} >{\raggedleft\arraybackslash}p{0.41\textwidth} @{}}
\textbf{\#} & \textbf{Proof} & \textbf{Mathematics} \\
\midrule
\textbf{1.} & We prove that join is above meet via transitivity for lattice on Order. & \\[0.45em]
\textbf{2.} & We split into exhaustive cases — the claim must hold in each one. & \\[0.45em]
\textbf{3.} & Case 1 (see step 2): suppose meet(a, b) <= b. & \\[0.45em]
\textbf{4.} & Case 2 (see step 2): suppose b <= join(a, b). & \\[0.45em]
\textbf{5.} & We invoke the derived fact governing meet on Order: meet is below the right argument, for meet on Order (instantiated for a, b). & \\[0.45em]
\textbf{6.} & We invoke the derived fact governing join on Order: join is above the right argument, for join on Order (instantiated for a, b). & \\[0.45em]
\textbf{7.} & We invoke the derived fact governing order on the natural numbers: less-or-equal is transitive, for order on the natural numbers (instantiated for meet(a, b), b, join(a, b)). & \\[0.45em]
\textbf{8.} & From step 4, step 5, step 6, and step 7, this implies meet(a, b) is at most join(a, b). Together with the other cases (step 3 and step 4), the goal is discharged. Hence proven. & $\displaystyle meet(a, b) \le join(a, b)$ \\[0.45em]
\bottomrule
\end{tabular}
\label{thm:Order-lattice-join_is_above_meet_via_transitivity}
\par\medskip\hrule
\end{document}
